% ---
% id: meat-core/prop:a-neutral
% kind: prop
% title: "Value neutrality along Alice's sectors"
% summary: >-
%   Pinching a feasible state along Alice's declared sectors keeps it feasible
%   with the same statistics and an objective no larger. So restricting the
%   variable to states block diagonal there changes no value, and the pinching
%   carries minimisers to minimisers.
% prominence: standard
% status: accepted
% uses: [def:declaration, def:feasible, def:homog, def:maps, def:branch, lem:branch-maps, thm:branch-formula, lem:cone-concavity, foundations/lem:kernel, foundations/def:pinching, def:admissible-declaration]
% mentions: [thm:reduction]
% assumes: []
% refs: []
% provenance: newly-derived
% lean:
% notes: "The proposition is independent of Step 1 of thm:reduction: either makes the restriction sound, and this one adds that nothing is lost."
% ---

Step~1 of the proof of Theorem~\ref{thm:reduction} shows that restricting the
variable to states block diagonal along Alice's declared sectors is sound; the
proposition below shows, independently, that the restriction also loses no
value.

\begin{proposition}
\label{prop:a-neutral}
Let $\Decl$ satisfy Definition~\ref{def:admissible-declaration} in the general format of
Definition~\ref{def:declaration}, in either source mode and with an arbitrary
list $\CList$, \tref{general} and \tref{party-local} entries included.
Let $\Feas^{\circ}$ be the set defined as $\Feas$ is in
Definition~\ref{def:feasible} but with block diagonality required along
$\SecSet$ only --- the feasible set of the same declaration with the factor
tagged $A$ forgotten --- and let
$\mathcal{P}_{\mathbb{T}}(X)\coloneqq\sum_{t}\Pi^{A}_tX\Pi^{A}_t$. Then for
every $\alpha\in(1,\infty)$, every distribution $\mathbf q$ on $\Chatset$ and
every $\omega\in\Feas^{\circ}$, the operator
$\omega'\coloneqq(\mathcal{P}_{\mathbb{T}}\otimes\idmap_{\Qsys})(\omega)$
satisfies
\begin{equation}\label{eq:a-neutral}
  \omega'\in\Feas,\qquad
  P_{\hat c}(\omega')=P_{\hat c}(\omega),\qquad
  K_{\hat c}(\omega')\ge K_{\hat c}(\omega)\quad\forall\hat c,\qquad
  O_{\alpha}(\mathbf q,\omega')\le O_{\alpha}(\mathbf q,\omega),
\end{equation}
the last in the extended-real sense. Consequently
$\inf_{\acSet\times\Feas}O_{\alpha}=\inf_{\acSet\times\Feas^{\circ}}O_{\alpha}$,
and $\omega\mapsto\omega'$ maps minimisers over $\Feas^{\circ}$ to minimisers
over $\Feas$.
\end{proposition}

\begin{proof}
$\mathcal{P}_{\mathbb{T}}$ is a pinching (\xref{foundations/def:pinching}):
completely positive, trace preserving, unital and self-adjoint for the trace
pairing.

\emph{Membership.} $\omega'\succeq0$ and $\tr\omega'=\tr\omega=1$.
$\mathcal{P}_{\mathbb{T}}\otimes\idmap_{\Qsys}$ commutes with
$X\mapsto(\idon{A}\otimes\Pi_s)X(\idon{A}\otimes\Pi_{s'})$, so $\omega'$ is
block diagonal along $\SecSet$ as $\omega$ is, and along $\mathbb{T}$ by
construction; hence along the joint sectors. In trusted-source mode
$\omega'_A=\mathcal{P}_{\mathbb{T}}(\omega_A)=\mathcal{P}_{\mathbb{T}}(\tau_A)
=\tau_A$, because $\tau_A$ is block diagonal along $\mathbb{T}$
(Definition~\ref{def:admissible-declaration}\,(iii)). For an entry read through its lifted
operator, the partial trace over $\Qsys$ commutes with a map acting on $A$
alone, so
$\tr_{\Qsys}[(\idon{A}\otimes\hat O_k)\omega']
=\mathcal{P}_{\mathbb{T}}\big(\tr_{\Qsys}[(\idon{A}\otimes\hat O_k)\omega]\big)$.
In trusted-source mode $\mathcal{P}_{\mathbb{T}}$ is positive and linear, so
$X\bowtie_k\Theta_k$ implies
$\mathcal{P}_{\mathbb{T}}(X)\bowtie_k\mathcal{P}_{\mathbb{T}}(\Theta_k)
=\Theta_k$, every $\Theta_k$ being block diagonal along $\mathbb{T}$: for a
\tref{homogeneous} entry $\Theta_k=c_k\tau_A$, for a general one by
its type (Definition~\ref{def:declaration}\,(DX2)). In untrusted-source mode
the entry value is the trace of the same operator, which
$\mathcal{P}_{\mathbb{T}}$ preserves. For a party-local entry
scoped to $f$, with $\hat D\coloneqq\hat O_k-c_k\idon{\Qsys_f}$, the partial
trace over $\Qsys_f$ commutes with a map acting on $A$ alone, so
$\tr_{\Qsys_f}[(\idon{}\otimes\hat D)\omega']
=(\mathcal{P}_{\mathbb{T}}\otimes\idmap)\big(
\tr_{\Qsys_f}[(\idon{}\otimes\hat D)\omega]\big)$ with the identity map on
the other measured factors; a pinching tensored with an identity map is a
pinching, hence positive, and $Y\bowtie_k0$ implies
$(\mathcal{P}_{\mathbb{T}}\otimes\idmap)(Y)\bowtie_k0$. So
$\omega'\in\Feas$.

\emph{Linear quantities.} For all $(x,v)$, by self-adjointness of the
pinching and block diagonality of $M_x$ along $\mathbb{T}$
(Definition~\ref{def:admissible-declaration}\,(iii)),
$\tr[(M_x\otimes F_v)\omega']
=\tr[(\mathcal{P}_{\mathbb{T}}(M_x)\otimes F_v)\omega]
=\tr[(M_x\otimes F_v)\omega]$. Hence $\Phimap[\omega']=\Phimap[\omega]$ by
\eqref{eq:Phimap}, so $P_{\hat c}$ is unchanged, and
$\tr\cG_{\hat c}(\omega')=\tr\cG_{\hat c}(\omega)$, that trace being
$\sum_{\kappa\in\Kann_{\mathrm{pass}}(\hat c)}\sum_{(x,v)\in a^{-1}(\kappa)}
\tr[(M_x\otimes F_v)\,\cdot\,]$ by the computation in the proof of
Lemma~\ref{lem:branch-maps}\,(i).

\emph{The non-linear term.} With $m\coloneqq|\mathbb{T}|$, an enumeration
$t\mapsto j_t\in\{1,\dots,m\}$ and $\zeta\coloneqq e^{2\pi i/m}$, the
operators $U_k\coloneqq\sum_t\zeta^{j_tk}\Pi^{A}_t$, $k=0,\dots,m-1$, are
unitaries on $A$ with
$\mathcal{P}_{\mathbb{T}}(X)=\frac1m\sum_kU_kXU_k^{\dagger}$, by the
unitary-mixture form of a pinching in \xref{foundations/def:pinching}, with
the projectors $\Pi^{A}_t$ enumerated by $j_t$. Each $U_k$ is a function of the
projectors $\Pi^{A}_t$ and therefore commutes with every operator block
diagonal along $\mathbb{T}$, in particular with every $\sqrt{M_x}$. Let
$\widetilde U_k$ be $U_k$ on the $A$ factor of the output space
$\Hs{R\,XY\,C_AC_B\,A\Qsys}$ of \eqref{eq:kraus} and the identity on the other
factors. Every term of $K_{\kappa}$ carries
$\sqrt{M_x}\otimes\sqrt{F_v}$ on $A\Qsys$, so
$K_{\kappa}(U_k\otimes\idon{\Qsys})=\widetilde U_kK_{\kappa}$ and, $\cG_{\hat
c}$ being linear,
\begin{equation*}
  \cG_{\hat c}(\omega')=\frac1m\sum_{k=0}^{m-1}
  \widetilde U_k\,\cG_{\hat c}(\omega)\,\widetilde U_k^{\dagger}.
\end{equation*}
For a unitary $U$ that is the identity on $R$ and any $X\succeq0$:
$(UXU^{\dagger})^{1/\alpha}=UX^{1/\alpha}U^{\dagger}$ by functional calculus;
$\cZ(UYU^{\dagger})=U\cZ(Y)U^{\dagger}$ because $U$ commutes with every
$Z_j=\dyad{j}_R$; and
$\tr[(U\cZ(Y)U^{\dagger})^{\alpha}]=\tr[\cZ(Y)^{\alpha}]$. So the functional
$T$ of \eqref{eq:Tcone} satisfies $T(UXU^{\dagger})=T(X)$, and by its
concavity on the whole cone (Lemma~\ref{lem:cone-concavity}\,(b))
\begin{equation*}
  T_{\hat c}(\omega')=T\big(\cG_{\hat c}(\omega')\big)
  \ge\frac1m\sum_kT\big(\widetilde U_k\cG_{\hat c}(\omega)
  \widetilde U_k^{\dagger}\big)=T_{\hat c}(\omega).
\end{equation*}
With \eqref{eq:fK},
$K_{\hat c}=T_{\hat c}+P_{\hat c}-\tr\cG_{\hat c}$, and the two linear terms
unchanged, $K_{\hat c}(\omega')\ge K_{\hat c}(\omega)$.

\emph{The objective and the infima.} In \eqref{eq:objective} the state
enters through $P_{\hat c}$, unchanged, and through
$\frac{1}{\alpha-1}d(q(\hat c),K_{\hat c})$ with $d$ the kernel of
\xref{foundations/lem:kernel}, non-increasing in its second argument; the
conventions at vanishing arguments are those of that kernel, so the
inequality holds in the extended-real sense. Finally
$\Feas\subseteq\Feas^{\circ}$, so the infimum over $\Feas$ is no smaller,
and by \eqref{eq:a-neutral} it is no larger; if $(\mathbf q,\omega)$ attains
the infimum over $\Feas^{\circ}$ then $(\mathbf q,\omega')$ attains the same
value over $\Feas$.
\end{proof}
